\documentclass[conference]{IEEEtran}
\usepackage[T1]{fontenc}
\usepackage{hyperref}
\hypersetup{colorlinks=true,linkcolor=black,urlcolor=blue}
\begin{document}

\title{Specialized LLM Skills Across Fixed Quantitative Portfolio Optimization Strategies}
\author{\IEEEauthorblockN{Carter Tierney}
\IEEEauthorblockA{\textit{Department of Computer Science}\\
\textit{Georgia State University}\\
Atlanta, GA, USA}}
\maketitle

\noindent{\small\bfseries\textit{Proposal Abstract:} I propose a controlled study of how large language model (LLM) views change the results of different quantitative portfolio optimization strategies. The data, the quant baseline, the Ledoit-Wolf covariance, the way views are built, and the guardrails are held fixed, so only the LLM's ratings and the strategy change. Each LLM rates every stock as a range (low, most likely, high) using a bare prompt, one of four specialized skills files, or a placebo (filler text the same length as a skills file). These ratings become Black-Litterman views that steer five fixed strategies, and each strategy is compared with its own quant-baseline-only version, which ranks stocks on momentum plus a pretrained forecasting model. I will track a weekly evaluation through the rank information coefficient (IC), portfolio returns and the Sharpe ratio over a live run to find which strategies the LLM helps or hurts the most, and why. To my knowledge, no prior study has tested LLM views across several fixed strategies on live data in the same period.\par}

\section{Research Questions}
\textit{How can LLM views built from specialized skills improve quantitative portfolio optimization strategies, and which strategies do LLMs help or hurt the most, and why?}

\noindent\textbf{RQ1} How often do the LLM's ratings alone beat the quant baseline's ratings?

\noindent\textbf{RQ2} How often does adding the LLM views improve each strategy's portfolio, compared with the same strategy run on the baseline alone (Q0)?

\noindent\textbf{RQ3} Do specialized skills change the result compared with bare prompts and placebo?

\noindent\textbf{RQ4} Which strategies do the LLM views help or hurt the most, and in which conditions (up or down markets, high or low volatility, earnings weeks)?

\noindent\textbf{RQ5} How much of the result comes from luck or general market movement rather than the LLM?

\section{Method: Weekly Pipeline from Input to Portfolio}
\begin{enumerate}
\item \textit{Data:} Prices, news and earnings dates for 100 S\&P 500 stocks balanced across sectors, frozen every Friday at market close. The number of stocks is set by the 1,000 GPU hours ARCTIC granted for the LLM ratings. Weekly moves of 10\% or more are kept and flagged, and a stock that falls 40\% or more is removed from the study.
\item \textit{Quant baseline score:} Built from price history alone, with no LLM. Each stock is ranked on its price momentum (12-1 month, 3-month and reversed 1-month returns) and on Kronos's next-week return forecast, then the two ranks are averaged. The spread of Kronos's 20 sampled forecast paths sets how much the baseline view is trusted.
\item \textit{LLM rating:} The model rates each stock as a range on a -5 to 5 scale (low, most likely, high) with a short rationale, using one prompt type: a bare prompt, one of four specialized skills files of similar length, or a placebo (filler text of the same length, with no investing content).
\item \textit{Residualize:} Remove the part of the LLM's most likely score that the baseline already explains, so the LLM view only adds new information.
\item \textit{Build views:} The most likely score and stock volatility set each view's expected return, centered on the baseline-only forecast so a neutral rating changes nothing, with a fixed view strength. The width of the range sets the uncertainty, so a narrow range like 3 to 5 counts more than a wide range like 1 to 5.
\item \textit{Black-Litterman:} Blend the baseline and LLM views with market-implied returns, using a Ledoit-Wolf covariance, into expected returns. The same expected returns and covariance go to every strategy.
\item \textit{Strategies:} Five fixed strategies turn the expected returns into a portfolio. Mean-variance maximizes return minus a risk penalty. Mean-CVaR maximizes return against the average loss in the worst 5\% of outcomes, using 10,000 Monte Carlo scenarios drawn from the Black-Litterman results. Robust mean-variance optimizes against the worst case within each view's uncertainty. Resampled mean-variance re-optimizes over 200 random draws of expected returns and averages the weights. Top-K holds the 20 stocks with the highest expected returns in equal weights.
\item \textit{Controls:} Each strategy also runs a control portfolio (Q0) using only the quant baseline view and no LLM views. For robust and resampled mean-variance, a confidence-only control (neutral rating, same view uncertainty) separates a view's information from its added confidence. Two more controls ignore forecasts entirely: hierarchical risk parity and an equal-weight portfolio.
\item \textit{Guardrails:} The same for every strategy: weekly rebalancing, daily checks with stop-loss -8\% and take-profit +15\%, a trading cost of 10 basis points (0.1\%) per trade, no shorts, and at most 5\% per stock.
\end{enumerate}

\section{Models}
Every model sees the same data and runs through the same pipeline. Open-weight models run on GSU ARCTIC through its OpenAI-compatible LiteLLM gateway, API models are called through their providers' batch APIs, and everything else (Kronos, the optimizers, the simulator and the statistics) runs on my own computer. All inputs are public market data and published news, and models banned by GSU's AI tool registry (such as Qwen and DeepSeek) are not used.
\begin{itemize}
\item \textit{Quant baseline model:} Kronos, a pretrained candlestick forecaster, run on my own GPU
\item \textit{Open-weight models:} Gemma 4 31B and MiniMax M2.7, already hosted on ARCTIC's gateway, so the study needs no new model deployments or gateway changes
\item \textit{API, cheap/small models:} two paid models, Claude Haiku 5.5 and GPT-6 Luna, using batch requests and prompt caching, and no frontier reasoning models
\end{itemize}

\section{Evaluation}
\textit{Scoring.} Each week, every LLM's most likely score is compared with the next week's returns using rank IC, the Spearman correlation between the ordering of scores and of returns. Rank IC is computed across stocks within one week, so a market-wide rise or fall does not change it. Portfolios are judged on their active return, the weekly return minus the return of the same strategy's Q0, so the strategy, costs and guardrails cancel out and only the LLM views differ. Every metric is saved for every portfolio, model, prompt and strategy, and null results are reported. Every test uses the saved ratings, weights and returns, so none of them reruns a model or an optimizer.

\textit{RQ1:} Compare the weekly rank IC of each LLM's ratings with that of the quant baseline, both computed on the same stocks each week. A residual IC, from a weekly Fama-MacBeth regression of next-week returns on the LLM score that controls for the baseline score, momentum, size, volatility, beta, sector, and each stock's trailing loadings on the value, profitability and investment factors, tests whether the LLM adds information beyond the baseline and known styles.

\textit{RQ2:} Compare the rank IC of the LLM-steered Black-Litterman tilt with that of the baseline-only (Q0) tilt, so the gap is zero when the LLM adds nothing. For each LLM portfolio, report the mean active return against Q0, the information ratio and tracking error, the Sharpe ratio gap tested with the Ledoit-Wolf bootstrap, the extra turnover, and the break-even trading cost at which the gain disappears. Results cover all weeks and, separately, only the weeks with an active LLM view.

\textit{RQ3:} Compare each skills file with the bare prompt and the placebo on both rank IC and active return, pooled across models in a regression with model fixed effects. Also test whether narrow ranges such as 3 to 5 predict next-week returns better than wide ranges such as 1 to 5.

\textit{RQ4:} Rank the strategies by how much the LLM views helped or hurt against their own Q0. Split weeks (market up or down, high or low volatility) and stocks (earnings week or not) into groups chosen in advance, and compare the gaps across groups. A strategy whose weights barely move with the views is reported as insensitive to views by design rather than as evidence that the LLM does not help. These results are exploratory.

\textit{RQ5:} Regress each portfolio's daily active return on the Fama-French five factors plus momentum, using Kenneth French's daily factor data. The loadings show which market or style exposure the LLM added, and the remaining alpha is the part that market movement and known styles cannot explain. The placebo prompt controls for prompt length and any change to the views, so a gain the placebo also shows does not come from the LLM's content, while chance is handled by the tests below.

\textit{Tests.} Weekly gaps use Newey-West standard errors with Student-t critical values, a block bootstrap 95\% confidence interval and a sign-flip permutation test. A small set of main comparisons, fixed before the live run, is corrected with Holm's method, while all other comparisons use the Benjamini-Hochberg false discovery rate and are reported as secondary. Each result also states the smallest effect the study could detect. Portfolio metrics include the annualized Sharpe ratio net of the Treasury bill rate (with its probabilistic version), the Sortino ratio, maximum drawdown and realized CVaR from daily returns, and paired comparisons with hierarchical risk parity and equal weight.

\section{Scope and Limitations}
A one-week pilot runs first, and then the live run starts on real weekly data. There is no backtest, since LLMs may have already seen news and prices from any past period. The live run covers one market period, so results may not hold if the market crashes or becomes highly volatile. Using 100 stocks, to fit ARCTIC's 1,000 GPU hours, makes small effects harder to detect. Adding strategies adds portfolios but not weeks, so with about 23 live weeks only moderate or larger gaps can be detected, which is why every result reports the smallest effect that the study could detect.

\section{Related Work}
The closest prior work puts LLM forecasts into Black-Litterman views, adds LLM signals to several portfolio methods, tests investor-style prompts and skills, and shows that backtests can leak future information into LLMs. This study holds the views, covariance and guardrails fixed and compares how the same LLM views change several optimization strategies, with placebo controls and a live run.

\section{Hypotheses}
\noindent\textbf{RQ1:} The LLM ratings alone will show a small positive rank IC that does not reliably beat the quant baseline, and most of it will disappear once known styles are controlled.

\noindent\textbf{RQ2:} Adding LLM views will raise the tilt's rank IC slightly over Q0, but weekly active returns will be small and mostly within noise, with the clearest effect in mean-variance and the smallest in Mean-CVaR and resampled mean-variance.

\noindent\textbf{RQ3:} Skills prompts will beat the bare prompt and placebo by a small margin, the placebo will perform about the same as the bare prompt, and narrow rating ranges such as 3 to 5 will predict next-week returns slightly better than wide ones.

\noindent\textbf{RQ4:} LLM views will help most in high-volatility and earnings weeks, when news moves prices most, and least in calm, low-volatility weeks with little news.

\noindent\textbf{RQ5:} Most of any return gap against Q0 will be explained by style exposure, especially to momentum and size, leaving an alpha close to zero once the factors are removed.

{\footnotesize\setlength{\itemsep}{0pt}
\begin{thebibliography}{51}
\setlength{\itemsep}{0pt}\setlength{\parskip}{0pt}
\bibitem{b1} Y. Lee, Y. Kim, J. Kim, S. Kim, and Y. Lee, ``\href{https://arxiv.org/abs/2504.14345}{LLM-Enhanced Black-Litterman Portfolio Optimization},'' arXiv:2504.14345, 2025.
\bibitem{b2} L. A. Mantshimuli and J. W. Muteba Mwamba, ``\href{https://link.springer.com/article/10.1007/s41060-026-01066-0}{Toward a Unified Agentic Framework for Regime-Aware Portfolio Optimization with LLM Signals},'' \textit{Int. J. Data Sci. Anal.}, vol. 22, Art. no. 92, 2026.
\bibitem{b3} Z. Li and L. Min, ``\href{https://www.iaeng.org/IJAM/issues_v55/issue_8/IJAM_55_8_21.pdf}{CVaR-based Portfolio with Investment Views Given by Large Language Models},'' \textit{IAENG Int. J. Appl. Math.}, vol. 55, no. 8, pp. 2512-2520, 2025.
\bibitem{b4} L. A. Mantshimuli and J. W. Muteba Mwamba, ``\href{https://www.mdpi.com/1911-8074/19/5/320}{Few-Shot Portfolio Optimization: Can Large Language Models Outperform Quantitative Portfolio Optimization?},'' \textit{J. Risk Financial Manag.}, vol. 19, no. 5, Art. no. 320, 2026.
\bibitem{b5} Y. Kim, Y. Lee, J. Kim, and Y. Lee, ``\href{https://arxiv.org/abs/2510.01664}{GuruAgents: Emulating Wise Investors with Prompt-Guided LLM Agents},'' arXiv:2510.01664, 2025.
\bibitem{b6} X. Li \textit{et al.}, ``\href{https://arxiv.org/abs/2602.12670}{SkillsBench: Benchmarking How Well Agent Skills Work Across Diverse Tasks},'' arXiv:2602.12670, 2026.
\bibitem{b7} A. Lopez-Lira and Y. Tang, ``\href{https://arxiv.org/abs/2304.07619}{Can ChatGPT Forecast Stock Price Movements? Return Predictability and Large Language Models},'' arXiv:2304.07619, 2023.
\bibitem{b8} C. Li \textit{et al.}, ``\href{https://arxiv.org/abs/2505.11065}{Time Travel is Cheating: Going Live with DeepFund for Real-Time Fund Investment Benchmarking},'' in \textit{Proc. NeurIPS Datasets and Benchmarks Track}, 2025.
\bibitem{b9} Z. Gao, W. Jiang, and Y. Yan, ``\href{https://arxiv.org/abs/2512.23847}{Detecting Lookahead Bias in LLM Forecasts},'' arXiv:2512.23847, 2025.
\bibitem{b10} Y. Shi \textit{et al.}, ``\href{https://arxiv.org/abs/2508.02739}{Kronos: A Foundation Model for the Language of Financial Markets},'' arXiv:2508.02739, 2025.
\bibitem{b11} F. Black and R. Litterman, ``\href{https://people.duke.edu/~charvey/Teaching/BA453_2006/Black_Litterman_Global_Portfolio_Optimization_1992.pdf}{Global Portfolio Optimization},'' \textit{Financial Analysts J.}, vol. 48, no. 5, pp. 28-43, 1992.
\bibitem{b12} O. Ledoit and M. Wolf, ``\href{https://econpapers.repec.org/RePEc:eee:jmvana:v:88:y:2004:i:2:p:365-411}{A Well-Conditioned Estimator for Large-Dimensional Covariance Matrices},'' \textit{J. Multivariate Anal.}, vol. 88, no. 2, pp. 365-411, 2004.
\bibitem{b13} R. T. Rockafellar and S. Uryasev, ``\href{https://www.risk.net/journal-risk/2161159/optimization-conditional-value-risk}{Optimization of Conditional Value-at-Risk},'' \textit{J. Risk}, vol. 2, no. 3, pp. 21-41, 2000.
\bibitem{b14} D. Goldfarb and G. Iyengar, ``\href{https://zhaosong-lu.github.io/ResearchPapers/robport-OMS.pdf}{Robust Portfolio Selection Problems},'' \textit{Math. Oper. Res.}, vol. 28, no. 1, pp. 1-38, 2003.
\bibitem{b15} R. Michaud and R. Michaud, ``\href{https://joim.com/articles-estimation-error-portfolio-optimization-resampling-solution}{Estimation Error and Portfolio Optimization: A Resampling Solution},'' \textit{J. Investment Manag.}, vol. 6, no. 1, 2008.
\bibitem{b16} M. L\'opez de Prado, ``\href{https://papers.ssrn.com/abstract=2708678}{Building Diversified Portfolios that Outperform Out-of-Sample},'' \textit{J. Portfolio Manag.}, vol. 42, no. 4, pp. 59-69, 2016.
\bibitem{b17} Y. Chen, B. T. Kelly, and D. Xiu, ``\href{https://papers.ssrn.com/sol3/papers.cfm?abstract_id=4416687}{Expected Returns and Large Language Models},'' SSRN Working Paper 4416687, 2022.
\bibitem{b18} S. He, L. Lv, A. Manela, and J. Wu, ``\href{https://arxiv.org/abs/2502.21206}{Chronologically Consistent Large Language Models},'' arXiv:2502.21206, 2025.
\bibitem{b19} Y. Abe, S. Matsuo, R. Kondo, and R. Hisano, ``\href{https://arxiv.org/abs/2411.19515}{Leveraging Large Language Models for Institutional Portfolio Management: Persona-Based Ensembles},'' arXiv:2411.19515, 2024.
\bibitem{b21} A. Voronina, O. Romanko, R. Cao, R. H. Kwon, and R. Mendoza-Arriaga, ``\href{https://arxiv.org/abs/2512.24526}{Generative AI-enhanced Sector-based Investment Portfolio Construction},'' arXiv:2512.24526, 2025.
\bibitem{b22} A. Kargarzadeh, N. Khaledian, N. Parvini, and A. Khaledian, ``\href{https://arxiv.org/abs/2608.12283}{Large Language Model-Driven Small-Capitalization Trading: Integrating Financial News Sentiment, Macroeconomic Indicators, and Technical Signals},'' arXiv:2608.12283, 2026.
\bibitem{b23} L. A. Mantshimuli and J. W. Muteba Mwamba, ``\href{https://doi.org/10.21511/imfi.22(3).2025.16}{Enhancing Portfolio Optimization with Multi-LLM Sentiment Aggregation: A Black-Litterman Integration Approach},'' \textit{Investment Manag. Financial Innov.}, vol. 22, no. 3, pp. 213-226, 2025.
\bibitem{b26} S. Wu \textit{et al.}, ``\href{https://arxiv.org/abs/2610.01213}{From Language-Model Stock Rankings to Testable Economic Rules: A Computational Audit},'' arXiv:2610.01213, 2026.
\bibitem{b27} Q. Duan \textit{et al.}, ``\href{https://arxiv.org/abs/2607.10286}{Can Agentic Trading Systems Pay for Their Own Intelligence?},'' arXiv:2607.10286, 2026.
\bibitem{b28} J. Chen and G. Wang, ``\href{https://arxiv.org/abs/2609.30705}{The Price of Thought: Does Test-Time Reasoning Pay in LLM Trading?},'' arXiv:2609.30705, 2026.
\bibitem{b29} B. Qu and M. Chen, ``\href{https://arxiv.org/abs/2606.29771}{CLQT: A Closed-Loop, Cost-Aware, Strategy-Consistent Benchmark for Diagnostic Evaluation of LLM Portfolio-Management Agents},'' arXiv:2606.29771, 2026.
\bibitem{b30} T. Zhu \textit{et al.}, ``\href{https://arxiv.org/abs/2605.28359}{From Knowing to Doing: A Memory-Controlled Benchmark for LLM Trading Agents on Stock Markets},'' arXiv:2605.28359, 2026.
\bibitem{b32} H. Yu, F. Li, and J. You, ``\href{https://arxiv.org/abs/2511.03628}{LiveTradeBench: Seeking Real-World Alpha with Large Language Models},'' arXiv:2511.03628, 2025.
\bibitem{b35} Y. Chen \textit{et al.}, ``\href{https://arxiv.org/abs/2510.02209}{StockBench: Can LLM Agents Trade Stocks Profitably In Real-world Markets?},'' arXiv:2510.02209, 2025.
\bibitem{b36} W. W. Li, H. Kim, M. Cucuringu, and T. Ma, ``\href{https://arxiv.org/abs/2505.07078}{Can LLM-based Financial Investing Strategies Outperform the Market in Long Run?},'' in \textit{Proc. KDD Datasets and Benchmarks Track}, 2026.
\bibitem{b37} Y. Ye \textit{et al.}, ``\href{https://arxiv.org/abs/2605.16895}{The Alpha Illusion: Reported Alpha from LLM Trading Agents Should Not Be Treated as Deployment Evidence},'' arXiv:2605.16895, 2026.
\bibitem{b38} J. Yao and Z. Zheng, ``\href{https://arxiv.org/abs/2606.08285}{Beyond Agent Architecture: Execution Assumptions and Reproducibility in LLM-Based Trading Systems},'' arXiv:2606.08285, 2026.
\bibitem{b39} S. S. Cao, W. Jiang, and H. Xu, ``\href{https://www.nber.org/papers/w35142}{Seeing the Goal, Missing the Truth: Human Accountability for AI Bias},'' NBER Working Paper 35142, 2026.
\bibitem{b41} H. Markowitz, ``\href{https://doi.org/10.1111/j.1540-6261.1952.tb01525.x}{Portfolio Selection},'' \textit{J. Finance}, vol. 7, no. 1, pp. 77-91, 1952.
\bibitem{b42} R. C. Grinold, ``\href{https://doi.org/10.3905/jpm.1989.409211}{The Fundamental Law of Active Management},'' \textit{J. Portfolio Manag.}, vol. 15, no. 3, pp. 30-37, 1989.
\bibitem{b43} E. F. Fama and K. R. French, ``\href{https://doi.org/10.1016/j.jfineco.2014.10.010}{A Five-Factor Asset Pricing Model},'' \textit{J. Financial Econ.}, vol. 116, no. 1, pp. 1-22, 2015.
\bibitem{b44} W. K. Newey and K. D. West, ``\href{https://doi.org/10.2307/1913610}{A Simple, Positive Semi-Definite, Heteroskedasticity and Autocorrelation Consistent Covariance Matrix},'' \textit{Econometrica}, vol. 55, no. 3, pp. 703-708, 1987.
\bibitem{b45} S. Holm, ``\href{https://www.jstor.org/stable/4615733}{A Simple Sequentially Rejective Multiple Test Procedure},'' \textit{Scand. J. Statist.}, vol. 6, no. 2, pp. 65-70, 1979.
\bibitem{b47} Y. Wang, T. G. You, L. Boussioux, and S. Liu, ``\href{https://arxiv.org/abs/2511.15202}{SOLID: a Framework of Synergizing Optimization and LLMs for Intelligent Decision-Making},'' arXiv:2511.15202, 2025.
\bibitem{b48} N. Deng, J. Wang, Q. Li, J. Zhang, and Y. Yang, ``\href{https://arxiv.org/abs/2609.24175}{FinRankGRPO: Optimizing LLMs for Listwise Financial Asset Ranking via Group Relative Policy Optimization},'' arXiv:2609.24175, 2026.
\bibitem{b49} Y. Zhao, S. Chen, and N. Su, ``\href{https://arxiv.org/abs/2605.27887}{PortBench: A Correlation-Aware, Full-Pipeline Benchmark for LLM-Driven Portfolio Management},'' arXiv:2605.27887, 2026.
\bibitem{b50} H. Cho and J. H. Kim, ``\href{https://arxiv.org/abs/2603.09301}{Constructing a Portfolio Optimization Benchmark Framework for Evaluating Large Language Models},'' arXiv:2603.09301, 2026.
\bibitem{b53} M. Florencio, ``\href{https://arxiv.org/abs/2607.20533}{Grounding Investor Views: Neural Predicates in the Black-Litterman Model},'' arXiv:2607.20533, 2026.
\bibitem{b54} W. Fu, ``\href{https://arxiv.org/abs/2510.05533}{The New Quant: A Survey of Large Language Models in Financial Prediction and Trading},'' arXiv:2510.05533, 2025.
\bibitem{b55} O. Ledoit and M. Wolf, ``\href{https://www.econ.uzh.ch/apps/workingpapers/wp/iewwp320.pdf}{Robust Performance Hypothesis Testing with the Sharpe Ratio},'' \textit{J. Empirical Finance}, vol. 15, no. 5, pp. 850-859, 2008.
\bibitem{b56} E. F. Fama and J. D. MacBeth, ``\href{https://doi.org/10.1086/260061}{Risk, Return, and Equilibrium: Empirical Tests},'' \textit{J. Political Econ.}, vol. 81, no. 3, pp. 607-636, 1973.
\bibitem{b57} Y. Benjamini and Y. Hochberg, ``\href{https://rss.onlinelibrary.wiley.com/doi/10.1111/j.2517-6161.1995.tb02031.x}{Controlling the False Discovery Rate: A Practical and Powerful Approach to Multiple Testing},'' \textit{J. Roy. Statist. Soc. B}, vol. 57, no. 1, pp. 289-300, 1995.
\bibitem{b58} M. A. Petersen, ``\href{https://papers.ssrn.com/sol3/papers.cfm?abstract_id=1320561}{Estimating Standard Errors in Finance Panel Data Sets: Comparing Approaches},'' \textit{Rev. Financial Stud.}, vol. 22, no. 1, pp. 435-480, 2009.
\bibitem{b59} V. DeMiguel, L. Garlappi, and R. Uppal, ``\href{https://econpapers.repec.org/RePEc:oup:rfinst:v:22:y:2009:i:5:p:1915-1953}{Optimal Versus Naive Diversification: How Inefficient is the 1/N Portfolio Strategy?},'' \textit{Rev. Financial Stud.}, vol. 22, no. 5, pp. 1915-1953, 2009.
\bibitem{b60} A. W. Lo, ``\href{https://www.tandfonline.com/doi/abs/10.2469/faj.v58.n4.2453}{The Statistics of Sharpe Ratios},'' \textit{Financial Analysts J.}, vol. 58, no. 4, pp. 36-52, 2002.
\bibitem{b61} K. R. French, ``\href{https://mba.tuck.dartmouth.edu/pages/faculty/ken.french/data_library.html}{Data Library},'' Tuck School of Business, Dartmouth College.
\end{thebibliography}}

\end{document}
